% practice_answers.tex -- Module 6 (Electromagnetism): worked answers to the Practice set.
% One practiceans per practiceq in practice.tex, in the same order (24 answers, 82 marks).
% Second, independent recalculation of every final answer:
% Directions were checked with vectors (F = q v x B, F = I l x B).
% MC answers keep their first lines together with \Needspace (conditional,
% layout-independent; no \clearpage or \pagebreak).
\begin{practiceanswers}

% ===================================================================== Section I
% ---- Q1 electron between plates
\Needspace{9\baselineskip}
\begin{practiceans}
  \correct{B}
  The field is uniform, so the electric force $F = qE$ has the same size and direction at
  every point. It acts at \ang{90} to the plates. Along the plates there is no force, so the
  velocity along the plates stays constant. Across the plates there is a constant
  acceleration. Constant velocity one way and constant acceleration at \ang{90} to it give a
  parabola. A ball thrown horizontally moves in the same way in the Earth's uniform
  gravitational field.
  \begin{itemize}
    \item \xmark\ A: ``uniform'' means the same everywhere, not zero. A constant force gives
          a constant acceleration, so the velocity changes.
    \item \xmark\ C: an arc of a circle needs a force that keeps turning to stay at
          \ang{90} to the velocity. That is the magnetic force. The electric force here
          always points the same way.
    \item \xmark\ D: the force has no part along the plates, so the velocity along the
          plates stays the same. The work done by the field increases only the velocity
          across the plates.
  \end{itemize}
\end{practiceans}

% ---- Q2 force on a wire
\Needspace{9\baselineskip}
\begin{practiceans}
  \correct{A}
  Use the right-hand palm rule. Point your fingers out of the page, towards you (the field).
  Point your thumb down the page (the current). Your palm faces to the left, so the force is
  to the left. (Vector check: $\vec{F} = I\vec{l} \times \vec{B}$ with $\vec{l}$ down and
  $\vec{B}$ out of the page gives a force to the left.)
  \begin{itemize}
    \item \xmark\ B: this is the answer you get with the left hand, or with the thumb
          pointing up the page.
    \item \xmark\ C: up the page is along the wire. The force is always at \ang{90} to the
          current.
    \item \xmark\ D: into the page is along the field. The force is always at \ang{90} to
          the field.
  \end{itemize}
\end{practiceans}

% ---- Q3 transformer current
\Needspace{9\baselineskip}
\begin{practiceans}
  \correct{C}
  $\dfrac{N_\mathrm{s}}{N_\mathrm{p}} = \dfrac{2400}{120} = 20$, so this transformer steps
  the voltage \emph{up} 20 times. An ideal transformer does not lose energy:
  $V_\mathrm{p}I_\mathrm{p} = V_\mathrm{s}I_\mathrm{s}$. So the secondary current is 20
  times \emph{smaller} than the primary current. This means the primary current is 20 times
  \emph{bigger} than the secondary current:
  $I_\mathrm{p} = I_\mathrm{s} \times \dfrac{N_\mathrm{s}}{N_\mathrm{p}}
  = \qty{0.25}{A} \times 20 = \qty{5.0}{A}$.
  \begin{itemize}
    \item \xmark\ A (\qty{0.0125}{A}): divides by 20. This uses the current ratio upside
          down. The power out would then be 400 times the power in.
    \item \xmark\ B (\qty{0.25}{A}): the current changes unless both coils have the same
          number of turns.
    \item \xmark\ D (\qty{100}{A}): multiplies by 20 twice ($20^2 = 400$).
  \end{itemize}
\end{practiceans}

% ---- Q4 Newton's third law
\Needspace{9\baselineskip}
\begin{practiceans}
  \correct{B}
  Currents in opposite directions repel. The force per metre on each wire is
  $\dfrac{F}{l} = \dfrac{\mu_0}{2\pi}\dfrac{I_1 I_2}{r}$. The same product $I_1 I_2$ is used
  for both wires, so the two forces have the same size. This is Newton's third law: the
  forces are an action--reaction pair, equal in size and opposite in direction. So the force
  on P also has size $F$.
  \begin{itemize}
    \item \xmark\ A: the forces are equal, but wires with opposite currents repel. They
          attract only when the currents are in the same direction.
    \item \xmark\ C: the bigger current in P does not make the force on P bigger. P carries
          3 times the current, but it sits in the weaker field made by Q. The two effects
          cancel.
    \item \xmark\ D: the force on P is not smaller either. Q sits in the stronger field
          made by P, but it carries only $\frac{1}{3}$ of the current. Again the effects
          cancel.
  \end{itemize}
\end{practiceans}

% ---- Q5 induction motor
\Needspace{9\baselineskip}
\begin{practiceans}
  \correct{C}
  The AC in the stator coils makes a rotating magnetic field. It sweeps past the bars of
  the rotor (the squirrel cage). The flux through the loops of the cage changes, so currents
  are induced in the bars (Faraday's law). These currents are in a magnetic field, so they
  feel forces (the motor effect). By Lenz's law, the forces oppose the \emph{relative}
  motion, so they drag the rotor round after the field. The rotor always turns a little
  more slowly than the field. This difference in speed is called the \textbf{slip}.
  \begin{itemize}
    \item \xmark\ A: nothing is wired to the rotor. Its currents are induced, so it needs no
          brushes and no commutator.
    \item \xmark\ B: at the same speed there is no relative motion, so the flux through the
          rotor does not change. Then no current is induced and there is no torque, so
          friction slows the rotor down.
    \item \xmark\ D: the rotor is a cage of copper or aluminium bars, not a magnet. Its
          magnetic field comes from the induced currents.
  \end{itemize}
\end{practiceans}

% ---- Q6 flux ratio
\Needspace{9\baselineskip}
\begin{practiceans}
  \correct{B}
  $\Phi = BA\cos\theta$, where $\theta$ is the angle between the field and the normal.
  The new area is $\pi(2r)^2 = 4\pi r^2$, so it is 4 times the old area. The new angle gives
  $\cos\ang{60} = 0.5$. So the new flux is $4 \times 0.5 \times \Phi = 2\Phi$.
  \begin{itemize}
    \item \xmark\ A ($\Phi$): doubles the area instead of multiplying it by 4 (forgets to
          square the radius): $2 \times 0.5 = 1$.
    \item \xmark\ C ($2\sqrt{3}\,\Phi \approx 3.5\Phi$): uses $\sin\ang{60}$ instead of
          $\cos\ang{60}$. The angle here is measured from the normal, so use $\cos$.
    \item \xmark\ D ($4\Phi$): forgets the angle.
  \end{itemize}
\end{practiceans}

% ---- Q7 jumping ring (Lenz)
\Needspace{9\baselineskip}
\begin{practiceans}
  \correct{A}
  Closing the switch makes the current in the solenoid grow from zero. An anticlockwise
  current (viewed from above) makes a field that points \emph{up} inside the solenoid
  (right-hand grip rule). So the upward flux through the ring \emph{increases}. By Lenz's
  law, the induced current in the ring opposes this increase. So it makes its own field
  that points \emph{down} inside the ring. A downward field, viewed from above, needs a
  clockwise current. Now the ring's current (clockwise) is opposite to the solenoid's
  current (anticlockwise). Opposite currents repel, so the ring is pushed up. (Pole view: the
  top of the solenoid is an N pole, and the bottom face of the ring is also an N pole. Like
  poles repel.) The push lasts only while the current grows. A steady current gives a steady
  flux, so no current is induced.
  \begin{itemize}
    \item \xmark\ B: the current direction is right, but opposite currents repel. A pull down
          would also move the ring into a stronger field. That would add to the change in
          flux instead of opposing it.
    \item \xmark\ C: an anticlockwise current would make an upward field. It would add to
          the growing flux. This breaks Lenz's law (and the conservation of energy).
    \item \xmark\ D: this is what happens when the switch is \emph{opened}. The upward flux
          then decreases, so the ring's current flows the same way as the solenoid's
          current (anticlockwise). Currents in the same direction attract.
  \end{itemize}
\end{practiceans}

% ---- Q8 proton and alpha particle
\Needspace{9\baselineskip}
\begin{practiceans}
  \correct{C}
  Speed after the voltage: $\frac{1}{2}mv^2 = qV$, so $v = \sqrt{\dfrac{2qV}{m}}$. Radius
  in the field: $r = \dfrac{mv}{qB}$. Put the speed into the radius:
  $r = \dfrac{m}{qB}\sqrt{\dfrac{2qV}{m}} = \dfrac{1}{B}\sqrt{\dfrac{2mV}{q}}$.
  $V$ and $B$ are the same for both particles, so $r$ is proportional to
  $\sqrt{\dfrac{m}{q}}$. For the alpha particle $\dfrac{m}{q}$ is
  $\dfrac{4}{2} = 2$ times the proton's value. So
  $\dfrac{r_\alpha}{r_\mathrm{p}} = \sqrt{2} \approx 1.41$.
  \begin{itemize}
    \item \xmark\ A ($\frac{1}{2}$): uses $\dfrac{q}{m}$ instead of $\dfrac{m}{q}$, and
          forgets the square root.
    \item \xmark\ B (1): the same voltage does not give the same speed or the same path.
    \item \xmark\ D (2): assumes both particles have the same speed, so that
          $r \propto \dfrac{m}{q}$. In fact the alpha particle is slower: its speed is
          $\dfrac{1}{\sqrt{2}}$ of the proton's speed.
  \end{itemize}
\end{practiceans}

% ===================================================================== Section II
% ---- Q9 electron gun
\begin{practiceans}
  \given{$V = \qty{2.5}{kV} = \qty{2500}{V}$; \ $d = \qty{1.0}{cm} = \qty{0.010}{m}$; \
    $q = \qty{1.602e-19}{C}$; \ $m_\mathrm{e} = \qty{9.109e-31}{kg}$; the electrons start
    from rest}
  \find{(a) the field $E$ \quad (b) the kinetic energy $K$ \quad (c) the speed $v$}
  \formula{$E = \dfrac{V}{d}$, \quad $W = qV$ (all the work becomes kinetic energy), \quad
    $K = \dfrac{1}{2}mv^2$}
  \substitute{\eqlines{E &= \frac{\qty{2500}{V}}{\qty{0.010}{m}} = \qty{2.5e5}{V.m^{-1}}\\
      K &= W = qV = \qty{1.602e-19}{C} \times \qty{2500}{V} = \qty{4.005e-16}{J}\\
      v^2 &= \frac{2K}{m_\mathrm{e}} \quad\text{(multiply both sides of
        $K = \tfrac{1}{2}mv^2$ by $\tfrac{2}{m}$)}\\
      &= \frac{2 \times \qty{4.005e-16}{J}}{\qty{9.109e-31}{kg}}
        = \qty{8.794e14}{m^2.s^{-2}}\\
      v &= \sqrt{\qty{8.794e14}{m^2.s^{-2}}} = \qty{2.965e7}{m.s^{-1}} \quad\text{(square root)}}}
  \answer{(a) \final{$E = \qty{2.5e5}{V.m^{-1}}$} \quad (b) \final{$K = \qty{4.0e-16}{J}$}
    \newline (c) \final{$v = \qty{3.0e7}{m.s^{-1}}$} (\qty{2.97e7}{m.s^{-1}})}
  \checkstep{$K = \qty{2500}{eV}$: the charge $e$ moved through \qty{2500}{V} \cmark. The
    speed is about 10\,\% of the speed of light: very fast, but well below $c$ \cmark. The
    gap width is not needed for (b) and (c): $W = qV$ does not depend on $d$.}
  \begin{markingcriteria}
    \criterion{(a) Correct field, \qty{2.5e5}{V.m^{-1}} (cm changed to m)}{1}
    \criterion{(b) Correct kinetic energy, \qty{4.0e-16}{J}}{1}
    \criterion{(c) Correct speed, \qty{3.0e7}{m.s^{-1}}}{1}
  \end{markingcriteria}
\end{practiceans}

% ---- Q10 busbars
\begin{practiceans}
  \given{$I_1 = I_2 = \qty{500}{A}$ (opposite directions); \ $r = \qty{4.0}{cm} =
    \qty{0.040}{m}$; \ $l = \qty{1.2}{m}$; \ $\mu_0 = 4\pi \times \qty{e-7}{N.A^{-2}}$}
  \find{(a) the force $F$ on each bar, and attract or repel \quad (b) $F/l$ for \qty{1}{A},
    \qty{1}{A} and \qty{1}{m}}
  \formula{$\dfrac{F}{l} = \dfrac{\mu_0}{2\pi}\dfrac{I_1 I_2}{r}$, \ so \
    $F = \dfrac{\mu_0}{2\pi}\dfrac{I_1 I_2}{r}\,l$}
  \substitute{\eqlines{\frac{\mu_0}{2\pi} &= \frac{4\pi \times \qty{e-7}{N.A^{-2}}}{2\pi}
        = \qty{2e-7}{N.A^{-2}} \quad\text{(cancel $2\pi$)}\\
      \frac{F}{l} &= \qty{2e-7}{N.A^{-2}} \times
        \frac{\qty{500}{A} \times \qty{500}{A}}{\qty{0.040}{m}} = \qty{1.25}{N.m^{-1}}\\
      F &= \qty{1.25}{N.m^{-1}} \times \qty{1.2}{m} = \qty{1.5}{N}
        \quad\text{(multiply by the length side by side)}}\par
    \textbf{(b)} $\dfrac{F}{l} = \qty{2e-7}{N.A^{-2}} \times
    \dfrac{\qty{1}{A} \times \qty{1}{A}}{\qty{1}{m}} = \qty{2e-7}{N.m^{-1}}$.}
  \answer{(a) \final{$F = \qty{1.5}{N}$} on each bar. The currents are opposite, so the bars
    \textbf{repel}. \quad (b) \final{\qty{2e-7}{N}} on each metre of each wire.}
  \checkstep{right-hand grip rule plus palm rule: each bar's field pushes the other bar
    \emph{away} \cmark. Forgetting to change \qty{4.0}{cm} to metres gives \qty{0.015}{N},
    100 times too small. Newton's third law: both bars feel the same \qty{1.5}{N}, in
    opposite directions.}
  \begin{markingcriteria}
    \criterion{(a) Correct force, \qty{1.5}{N} ($r$ in metres)}{1}
    \criterion{(a) States that the bars repel}{1}
    \criterion{(b) \qty{2e-7}{N} per metre}{1}
  \end{markingcriteria}
\end{practiceans}

% ---- Q11 neon sign transformer
\begin{practiceans}
  \given{$V_\mathrm{p} = \qty{240}{V}$; \ $V_\mathrm{s} = \qty{7.2}{kV} = \qty{7200}{V}$;
    \ $N_\mathrm{p} = 300$; \ $I_\mathrm{s} = \qty{20}{mA} = \qty{0.020}{A}$; \ real
    $I_\mathrm{p} = \qty{0.75}{A}$}
  \find{(a) $N_\mathrm{s}$ \quad (b) ideal $I_\mathrm{p}$ \quad (c) the efficiency
    $\eta$}
  \formula{$\dfrac{V_\mathrm{p}}{V_\mathrm{s}} = \dfrac{N_\mathrm{p}}{N_\mathrm{s}}$, \quad
    $V_\mathrm{p}I_\mathrm{p} = V_\mathrm{s}I_\mathrm{s}$, \quad
    $\eta = \dfrac{P_\mathrm{out}}{P_\mathrm{in}} \times 100\,\%$ with $P = VI$}
  \substitute{\eqlines{N_\mathrm{s} &= N_\mathrm{p} \times \frac{V_\mathrm{s}}{V_\mathrm{p}}
        = 300 \times \frac{\qty{7200}{V}}{\qty{240}{V}} = 300 \times 30 = \num{9000}\\
      I_\mathrm{p} &= \frac{V_\mathrm{s}I_\mathrm{s}}{V_\mathrm{p}}
        = \frac{\qty{7200}{V} \times \qty{0.020}{A}}{\qty{240}{V}}
        = \frac{\qty{144}{W}}{\qty{240}{V}} = \qty{0.60}{A}\\
      \eta &= \frac{\qty{144}{W}}{\qty{240}{V} \times \qty{0.75}{A}} \times 100\,\%
        = \frac{\qty{144}{W}}{\qty{180}{W}} \times 100\,\% = 80\,\%}}
  \answer{(a) \final{$N_\mathrm{s} = \num{9000}$ turns} \quad
    (b) \final{$I_\mathrm{p} = \qty{0.60}{A}$} \quad (c) \final{$\eta = 80\,\%$}}
  \checkstep{a step-up transformer has more turns on the secondary \cmark. The primary
    current is bigger than the secondary current \cmark. The real primary current
    (\qty{0.75}{A}) is bigger than the ideal one (\qty{0.60}{A}), because some energy is
    wasted \cmark.}
  \begin{markingcriteria}
    \criterion{(a) \num{9000} turns}{1}
    \criterion{(b) \qty{0.60}{A}, with mA changed to A}{1}
    \criterion{(c) Efficiency 80\,\%}{1}
  \end{markingcriteria}
\end{practiceans}

% ---- Q12 hanging rod
\begin{practiceans}
  \given{$l = \qty{0.20}{m}$; \ $m = \qty{8.0}{g} = \qty{0.0080}{kg}$; \
    $B = \qty{0.050}{T}$ into the page; \ $\theta = \ang{90}$; \ $g = \qty{9.8}{m.s^{-2}}$}
  \find{(a) the current $I$ \quad (b) its direction \quad (c) the current when
    $\theta = \ang{30}$}
  \formula{tension zero means the magnetic force holds up the weight:
    $lIB\sin\theta = mg$, \ so \ $I = \dfrac{mg}{lB\sin\theta}$}
  \substitute{\eqlines{mg &= \qty{0.0080}{kg} \times \qty{9.8}{m.s^{-2}} = \qty{0.0784}{N}\\
      I &= \frac{\qty{0.0784}{N}}{\qty{0.20}{m} \times \qty{0.050}{T} \times \sin\ang{90}}
        = \frac{\qty{0.0784}{N}}{\qty{0.010}{T.m}} = \qty{7.84}{A}}\par
    \textbf{(b)} The force must point up. Palm rule: fingers into the page (field), palm
    facing up (force). Your thumb then points to the right, so the current flows to the
    right.\par
    \textbf{(c)} Both the rod and the field are still horizontal, so the force is still
    vertical. Only $\sin\theta$ changes:\par\nobreak
    \eqlines{I &= \frac{\qty{0.0784}{N}}{\qty{0.20}{m} \times \qty{0.050}{T} \times \sin\ang{30}}
      = \frac{\qty{0.0784}{N}}{\qty{0.0050}{T.m}} = \qty{15.68}{A}}}
  \answer{(a) \final{$I = \qty{7.8}{A}$} \quad (b) to the \textbf{right} \quad
    (c) \final{$I = \qty{16}{A}$}}
  \checkstep{at \ang{30} only half of the field is perpendicular to the rod
    ($\sin\ang{30} = 0.5$), so twice the current is needed \cmark. Using $\cos\ang{30}$ by
    mistake gives \qty{9.1}{A}.}
  \begin{markingcriteria}
    \criterion{(a) Equates the magnetic force and the weight, $lIB = mg$ (g changed to kg)}{1}
    \criterion{(a) Correct current, \qty{7.8}{A}}{1}
    \criterion{(b) To the right, with a correct use of the right-hand palm rule}{1}
    \criterion{(c) Correct current, \qty{16}{A}, using $\sin\ang{30}$}{1}
  \end{markingcriteria}
\end{practiceans}

% ---- Q13 flipped coil
\begin{practiceans}
  \given{$N = 75$; \ diameter \qty{4.0}{cm}, so $r = \qty{2.0}{cm} = \qty{0.020}{m}$; \
    $B = \qty{0.35}{T}$; \ $\theta$ goes from \ang{0} to \ang{180} (the normal turns round);
    \ $\Delta t = \qty{0.060}{s}$}
  \find{(a) $\Phi_\mathrm{initial}$ \quad (b) the size of the average emf \quad (c) two other
    ways}
  \formula{$\Phi = BA\cos\theta$ with $A = \pi r^2$, \quad
    $\varepsilon = -N\dfrac{\Delta\Phi}{\Delta t}$}
  \substitute{\eqlines{A &= \pi \times (\qty{0.020}{m})^2 = \qty{1.257e-3}{m^2}\\
      \Phi_\mathrm{initial} &= \qty{0.35}{T} \times \qty{1.257e-3}{m^2} \times \cos\ang{0}
        = \qty{4.398e-4}{Wb}\\
      \Phi_\mathrm{final} &= \qty{0.35}{T} \times \qty{1.257e-3}{m^2} \times \cos\ang{180}
        = \qty{-4.398e-4}{Wb}\\
      \Delta\Phi &= \Phi_\mathrm{final} - \Phi_\mathrm{initial} = \qty{-8.796e-4}{Wb}\\
      \varepsilon &= -75 \times \frac{\qty{-8.796e-4}{Wb}}{\qty{0.060}{s}}
        = \qty{1.100}{V}}\par
    \textbf{(c)} Change the field strength $B$: move the magnet closer or further away, or
    change the current in an electromagnet. Or change the area $A$ inside the field: squash
    the coil, or move it into or out of the field.}
  \answer{(a) \final{$\Phi = \qty{4.4e-4}{Wb}$} \quad
    (b) \final{$\varepsilon = \qty{1.1}{V}$} \quad (c) change $B$; change $A$}
  \checkstep{after the flip, the same field lines go through the coil from the other side.
    So the flux changes by $2BA$, not by 0 \cmark. Common mistakes: using the diameter as
    the radius gives \qty{4.4}{V}; forgetting the factor 2 gives \qty{0.55}{V}.}
  \begin{markingcriteria}
    \criterion{(a) Correct flux, \qty{4.4e-4}{Wb} (radius, not diameter)}{1}
    \criterion{(b) Shows that $\Delta\Phi = 2BA = \qty{8.8e-4}{Wb}$}{1}
    \criterion{(b) Correct emf, \qty{1.1}{V}}{1}
    \criterion{(c) Two correct ways (change $B$; change the area in the field)}{1}
  \end{markingcriteria}
\end{practiceans}

% ---- Q14 torque on a motor coil
\begin{practiceans}
  \given{$n = 120$; \ $A = \qty{0.030}{m} \times \qty{0.050}{m} = \qty{1.5e-3}{m^2}$; \
    $I = \qty{0.75}{A}$; \ $B = \qty{0.18}{T}$}
  \find{(a) $\tau_\mathrm{max}$ \quad (b) $\tau$ when the plane is at \ang{20} to $B$ \quad
    (c) the job of the commutator}
  \formula{$\tau = nIAB\sin\theta$, where $\theta$ is the angle between $B$ and the
    \emph{normal} to the coil}
  \substitute{\textbf{(a)} The torque is largest when $\theta = \ang{90}$ (plane of the coil
    parallel to $B$):\par\nobreak
    \eqlines{\tau_\mathrm{max} &= 120 \times \qty{0.75}{A} \times \qty{1.5e-3}{m^2} \times
      \qty{0.18}{T} \times \sin\ang{90} = \qty{0.0243}{N.m}}\par
    \textbf{(b)} The normal is at \ang{90} to the plane. So if the plane makes \ang{20} with
    $B$, the normal makes $\ang{90} - \ang{20} = \ang{70}$ with $B$:\par\nobreak
    \eqlines{\tau &= \qty{0.0243}{N.m} \times \sin\ang{70} = \qty{0.0243}{N.m} \times 0.9397
      = \qty{0.02283}{N.m}}\par
    \textbf{(c)} Every half turn, the coil passes the position where its plane is at
    \ang{90} to the field. After that, the forces on its sides would turn it back the other
    way. The split-ring commutator reverses the current in the coil at this moment. So the
    forces, and the torque, keep turning the coil the same way.}
  \answer{(a) \final{$\tau_\mathrm{max} = \qty{0.024}{N.m}$} \quad
    (b) \final{$\tau = \qty{0.023}{N.m}$}}
  \checkstep{the plane is close to parallel with $B$, so the torque is close to the maximum
    \cmark. Using $\sin\ang{20}$ by mistake gives \qty{0.0083}{N.m}. Second method (taking the
    \qty{5.0}{cm} sides parallel to the axle): each of these sides feels
    $nlIB = 120 \times 0.050 \times 0.75 \times 0.18 = \qty{0.81}{N}$. Its lever arm is
    $\qty{0.015}{m} \times \cos\ang{20}$. The two sides give \qty{0.0228}{N.m} \cmark.}
  \begin{markingcriteria}
    \criterion{(a) Correct maximum torque, \qty{0.024}{N.m} (cm changed to m)}{1}
    \criterion{(b) Uses $\theta = \ang{70}$ (the angle between $B$ and the normal)}{1}
    \criterion{(b) Correct torque, \qty{0.023}{N.m}}{1}
    \criterion{(c) The commutator reverses the current every half turn, so the torque keeps
      the same direction}{1}
  \end{markingcriteria}
\end{practiceans}

% ---- Q15 zipline brake
\begin{practiceans}
  \textbf{(a)}
  \given{$m = \qty{90}{kg}$; \ $u = \qty{15}{m.s^{-1}}$; \ $v = \qty{3.0}{m.s^{-1}}$; no
    friction, no air resistance; the track is flat}
  \find{the heat made by the brake}
  \formula{the kinetic energy lost becomes heat: heat $= \dfrac{1}{2}mu^2 - \dfrac{1}{2}mv^2$}
  \substitute{\eqlines{\tfrac{1}{2}mu^2 &= \tfrac{1}{2} \times \qty{90}{kg} \times
        (\qty{15}{m.s^{-1}})^2 = \qty{10125}{J}\\
      \tfrac{1}{2}mv^2 &= \tfrac{1}{2} \times \qty{90}{kg} \times (\qty{3.0}{m.s^{-1}})^2
        = \qty{405}{J}\\
      \text{heat} &= \qty{10125}{J} - \qty{405}{J} = \qty{9720}{J}}}
  \answer{\final{heat $= \qty{9.7e3}{J}$}}
  \checkstep{using $\frac{1}{2}m(u - v)^2$ is wrong: it gives \qty{6480}{J}. Find each
    kinetic energy first, then subtract. The brake removes 96\,\% of the kinetic energy.}

  \begin{modelanswer}
  \textbf{(b)} As the fin moves through the field, the flux through loops in the fin changes. This
  induces emfs and eddy currents in the fin (Faraday's law,
  $\varepsilon = -N\dfrac{\Delta\Phi}{\Delta t}$). By Lenz's law, the eddy currents oppose
  the motion: the field pushes on them with a backward force ($F = lIB$).

  How fast the flux changes depends on the speed of the fin. When the trolley slows down,
  the flux changes more slowly. So the emf is smaller, the eddy currents are smaller, and
  the backward force on them is smaller. That is why the trolley slows quickly at first, then
  more gently. (So a magnetic brake alone cannot bring the trolley fully to rest.)
  \end{modelanswer}
  \begin{markingcriteria}
    \criterion{(a) Correct heat, \qty{9.7e3}{J}}{1}
    \criterion{(b) Lower speed means a slower change of flux, so a smaller induced emf and
      smaller eddy currents}{1}
    \criterion{(b) Smaller eddy currents give a smaller opposing (Lenz) force}{1}
  \end{markingcriteria}
\end{practiceans}

% ---- Q16 transmission
\begin{practiceans}
  \given{$P = \qty{18}{MW} = \qty{1.8e7}{W}$; \ $R = \qty{4.2}{\ohm}$; \
    $V = \qty{66}{kV} = \qty{6.6e4}{V}$, then $V = \qty{330}{kV} = \qty{3.3e5}{V}$}
  \find{(a) $I$ and $P_\mathrm{loss}$ at \qty{66}{kV} \quad (b) $P_\mathrm{loss}$ at
    \qty{330}{kV} \quad (c) the reason for high voltage}
  \formula{$P = VI$, so $I = \dfrac{P}{V}$; \quad $P_\mathrm{loss} = I^2R$ (not on the sheet:
    from $P = VI$ and $V = IR$)}
  \substitute{\eqlines{I &= \frac{\qty{1.8e7}{W}}{\qty{6.6e4}{V}} = \qty{272.7}{A}\\
      P_\mathrm{loss} &= (\qty{272.7}{A})^2 \times \qty{4.2}{\ohm} = \qty{3.124e5}{W}\\
      I &= \frac{\qty{1.8e7}{W}}{\qty{3.3e5}{V}} = \qty{54.55}{A}
        \quad\text{(at \qty{330}{kV})}\\
      P_\mathrm{loss} &= (\qty{54.55}{A})^2 \times \qty{4.2}{\ohm} = \qty{12498}{W}}}
  \answer{(a) \final{$I = \qty{2.7e2}{A}$}, \ \final{$P_\mathrm{loss} = \qty{3.1e5}{W}$}
    (\qty{312}{kW}) \newline (b) \final{$P_\mathrm{loss} = \qty{1.2e4}{W}$}
    (\qty{12498}{W} rounded to 2 significant figures)}
  \checkstep{\textbf{(c)} The power sent is the same, so 5 times the voltage gives
    $\frac{1}{5}$ of the current ($P = VI$). The loss depends on the \emph{square} of the
    current ($I^2R$), so it falls to $\frac{1}{25}$:
    $\dfrac{\qty{312400}{W}}{\qty{12500}{W}} = 25$ \cmark. The loss drops from 1.7\,\% to
    0.07\,\% of the power sent. Less energy is wasted as heat in the lines. \textbf{Trap:}
    $V^2/R$ with $V = \qty{66}{kV}$ gives \qty{1.04e9}{W}, more than the power sent. In
    $V^2/R$, $V$ must be the voltage \emph{drop} along the line, not the transmission
    voltage.}
  \begin{markingcriteria}
    \criterion{(a) Correct current, \qty{2.7e2}{A}}{1}
    \criterion{(a) Correct power loss, \qty{3.1e5}{W}}{1}
    \criterion{(b) Correct power loss, \qty{1.2e4}{W} (\qty{1.25e4}{W} is also correct)}{1}
    \criterion{(c) Higher voltage gives a smaller current for the same power, and the loss
      $I^2R$ falls with the square of the current (here 25 times smaller)}{1}
  \end{markingcriteria}
\end{practiceans}

% ---- Q17 flux-time graph
\begin{practiceans}
  \given{$N = 30$. From the graph: \qtyrange{0}{40}{ms}: $\Phi$ rises from 0 to
    \qty{1.2}{mWb}; \qtyrange{40}{60}{ms}: steady; \qtyrange{60}{80}{ms}: falls from
    \qty{1.2}{mWb} to \qty{-1.2}{mWb}; \qtyrange{80}{100}{ms}: steady}
  \find{the emf in each part}
  \formula{$\varepsilon = -N\dfrac{\Delta\Phi}{\Delta t}$, where
    $\dfrac{\Delta\Phi}{\Delta t}$ is the gradient of the graph}
  \substitute{\eqlines{\text{0--40 ms:}\ \frac{\Delta\Phi}{\Delta t}
        &= \frac{\qty{1.2e-3}{Wb} - 0}{\qty{0.040}{s}} = \qty{0.030}{Wb.s^{-1}}\\
      \varepsilon &= -30 \times \qty{0.030}{Wb.s^{-1}} = \qty{-0.90}{V}\\
      \text{40--60 ms:}\ \frac{\Delta\Phi}{\Delta t} &= 0, \ \text{so}\ \varepsilon = 0}\par
    \eqlines{\text{60--80 ms:}\ \frac{\Delta\Phi}{\Delta t}
        &= \frac{\qty{-1.2e-3}{Wb} - \qty{1.2e-3}{Wb}}{\qty{0.020}{s}}
        = \qty{-0.12}{Wb.s^{-1}}\\
      \varepsilon &= -30 \times (\qty{-0.12}{Wb.s^{-1}}) = +\qty{3.6}{V}\\
      \text{80--100 ms:}\ \frac{\Delta\Phi}{\Delta t} &= 0, \ \text{so}\ \varepsilon = 0}}
  \answer{(a) \final{$\lvert\varepsilon\rvert = \qty{0.90}{V}$} \quad (c) see the graph:
    \qty{-0.90}{V}, then 0, then $+\qty{3.6}{V}$, then 0}
  \checkstep{\textbf{(b)} The emf depends on how fast the flux \emph{changes} (the gradient),
    not on the size of the flux. At \qty{70}{ms} the flux is zero, but it is changing at its
    fastest rate (\qty{0.12}{Wb.s^{-1}}). So the emf is \qty{3.6}{V}, not zero.
    \textbf{Area check:} $\qty{0.90}{V} \times \qty{0.040}{s} = \qty{0.036}{V.s}
    = N\lvert\Delta\Phi\rvert = 30 \times \qty{1.2e-3}{Wb}$ \cmark.}
  \begin{figure}
    \centering
    \begin{tikzpicture}
      \begin{axis}[hsc, width=0.78\linewidth, height=0.34\linewidth,
          xmin=0, xmax=100, ymin=-4, ymax=4, xlabel={$t$ (\si{ms})},
          ylabel={$\varepsilon$ (\si{V})}, xtick={0,10,...,100}, ytick={-4,-3,...,4}]
        \draw[hscink!70, line width=0.6pt] (axis cs:0,0) -- (axis cs:100,0);
        \addplot[hscseriesB, solid, line width=1.3pt, mark=none, const plot]
          coordinates {(0,-0.9) (40,0) (60,3.6) (80,0) (100,0)};
        \node[anchor=north, font=\footnotesize] at (axis cs:20,-1.0) {$-0.90$};
        \node[anchor=south, font=\footnotesize] at (axis cs:70,3.65) {$+3.6$};
      \end{axis}
    \end{tikzpicture}
    \caption*{Answer to \cref{pq:fluxgraph}(c): the emf induced in coil S.}
  \end{figure}
  \begin{markingcriteria}
    \criterion{(a) Correct emf, \qty{0.90}{V}}{1}
    \criterion{(b) The emf depends on the rate of change of flux (the gradient), which is
      largest at \qty{70}{ms}}{1}
    \criterion{(c) Correct shape: constant blocks, with zero emf where the flux is steady}{1}
    \criterion{(c) Correct values: \qty{-0.90}{V} and $+\qty{3.6}{V}$ (opposite signs, the
      second 4 times as big)}{1}
  \end{markingcriteria}
\end{practiceans}

% ---- Q18 mass spectrometer
\begin{practiceans}
  \textbf{(a)} The magnetic force is always at \ang{90} to the velocity, so it is the
  centripetal force:
  \begin{align*}
    qvB\sin\ang{90} &= \frac{mv^2}{r} && \text{magnetic force = centripetal force}\\
    qB &= \frac{mv}{r} && \text{divide both sides by $v$ ($\sin\ang{90} = 1$)}\\
    r &= \frac{mv}{qB} && \text{multiply both sides by $\dfrac{r}{qB}$}
  \end{align*}

  \textbf{(b)}
  \given{$q = \qty{1.602e-19}{C}$; \ $v = \qty{2.0e5}{m.s^{-1}}$; \ $B = \qty{0.40}{T}$; \
    $m_6 = \qty{9.99e-27}{kg}$; \ $m_7 = \qty{1.165e-26}{kg}$}
  \find{the distance PQ}
  \formula{$r = \dfrac{mv}{qB}$. Each ion moves through a semicircle, so it lands a
    \emph{diameter} ($2r$) from the entrance.}
  \substitute{\eqlines{qB &= \qty{1.602e-19}{C} \times \qty{0.40}{T} = \num{6.408e-20}\\
      r_6 &= \frac{\qty{9.99e-27}{kg} \times \qty{2.0e5}{m.s^{-1}}}{\num{6.408e-20}}
        = \qty{0.03118}{m}\\
      r_7 &= \frac{\qty{1.165e-26}{kg} \times \qty{2.0e5}{m.s^{-1}}}{\num{6.408e-20}}
        = \qty{0.03636}{m}\\
      \text{PQ} &= 2r_7 - 2r_6 = \qty{0.07272}{m} - \qty{0.06236}{m} = \qty{0.01036}{m}}}
  \answer{\final{PQ $= \qty{1.0e-2}{m}$} (\qty{1.0}{cm}). The lighter lithium-6 ions land at
    P.}
  \checkstep{$r \propto m$ (same $v$, $q$ and $B$): $\dfrac{r_7}{r_6} = 1.166 =
    \dfrac{m_7}{m_6}$ \cmark. Forgetting that the ions land a \emph{diameter} away gives
    \qty{5.2}{mm}. The ions are positive, $v$ is up the page and $B$ is out of the page.
    So the palm rule gives a force to the right at the start. The ions curve to the right,
    as drawn \cmark.}
  \begin{markingcriteria}
    \criterion{(a) Sets the magnetic force equal to the centripetal force,
      $qvB = \dfrac{mv^2}{r}$}{1}
    \criterion{(a) Rearranges correctly to $r = \dfrac{mv}{qB}$}{1}
    \criterion{(b) Correct radii, \qty{3.12}{cm} and \qty{3.64}{cm}}{1}
    \criterion{(b) Uses diameters: PQ $= \qty{1.0}{cm}$}{1}
  \end{markingcriteria}
\end{practiceans}

% ===================================================================== Section III
% ---- Q19 evaluate the ideal-transformer claim
\begin{practiceans}
  \begin{modelanswer}
  \textbf{Using the data.} An ideal 12\,:\,1 transformer would give
  $V_\mathrm{s} = \dfrac{\qty{240}{V}}{12} = \qty{20.0}{V}$. The measured value is
  \qty{18.6}{V}, only 93\,\% of the ideal value.
  \begin{itemize}
    \item Power in: $P_\mathrm{in} = V_\mathrm{p}I_\mathrm{p} = \qty{240}{V} \times
          \qty{0.275}{A} = \qty{66.0}{W}$.
    \item Power out: $P_\mathrm{out} = V_\mathrm{s}I_\mathrm{s} = \qty{18.6}{V} \times
          \qty{3.20}{A} = \qty{59.5}{W}$.
    \item Efficiency: $\eta = \dfrac{\qty{59.5}{W}}{\qty{66.0}{W}} \times 100\,\% = 90\,\%$.
          So \qty{6.5}{W} (about 10\,\%) is wasted, mostly as heat.
    \item Current ratio: $\dfrac{I_\mathrm{s}}{I_\mathrm{p}} = \dfrac{\qty{3.20}{A}}{\qty{0.275}{A}}
          = 11.6$. An ideal transformer would give exactly 12.
  \end{itemize}

  \textbf{Causes of the loss, and how designers reduce them.}
  \begin{itemize}
    \item \textbf{Resistive heating.} The copper coils have resistance, so the current heats
          them ($P = I^2R$). The secondary carries the larger current (\qty{3.20}{A}), so
          its heating matters most. Its resistance also lowers the output voltage.
          \emph{Fix:} use thick, low-resistance copper wire, especially for the
          high-current coil.
    \item \textbf{Eddy currents.} The changing flux induces currents in the iron core,
          which heat it. \emph{Fix:} a laminated core (thin iron sheets, insulated from each
          other) stops large eddy currents.
    \item \textbf{Incomplete flux linkage.} Some flux from the primary leaks into the air
          and misses the secondary. So less emf is induced, and the output voltage is lower
          than the turns ratio predicts. \emph{Fix:} wind both coils on one closed iron
          core, one on top of the other.
    \item (Also: the core is magnetised back and forth 50 times a second, which wastes a
          little energy. \emph{Fix:} a soft-iron core.)
  \end{itemize}

  \textbf{Judgement.} The claim is not correct, but the ideal model is a useful first
  approximation. This transformer is quite efficient (90\,\%), and the ideal equations
  predict the output voltage to within about 7\,\%. But about 10\,\% of the input energy is
  wasted, and the output voltage is lower than the turns ratio predicts. So the transformer
  does not behave exactly like an ideal one.
  \end{modelanswer}
  \begin{markingcriteria}
    \criterion{Makes a clear judgement supported by the data (efficiency AND the ideal
      output voltage); explains at least two causes of energy loss AND links each to a design
      strategy that reduces it}{5}
    \criterion{Makes a judgement using the efficiency or the voltage; explains two causes
      of loss and gives strategies}{4}
    \criterion{Calculates the efficiency or compares the voltages, AND describes two causes
      of loss (or one cause with a strategy)}{3}
    \criterion{Calculates the input and output power, OR identifies two causes of loss or
      two strategies}{2}
    \criterion{Identifies one relevant cause of energy loss or one strategy}{1}
  \end{markingcriteria}
\end{practiceans}

% ---- Q20 generator
\begin{practiceans}
  \textbf{(a)} One turn takes $T = \dfrac{1}{25} = \qty{0.040}{s}$, so two turns take
  \qty{0.080}{s}. At the start, the plane of the coil is at right angles to the field. The
  flux is at its largest but not changing for an instant, so the emf is zero. The output is
  a sine curve that starts at zero, with a peak of \qty{4.0}{V} and a period of
  \qty{0.040}{s}. (Starting downwards is also correct: it depends on which way the coil
  turns.)
  \begin{figure}
    \centering
    \begin{tikzpicture}
      \begin{axis}[hsc, width=0.78\linewidth, height=0.34\linewidth,
          xmin=0, xmax=0.08, ymin=-5, ymax=5, xlabel={$t$ (\si{s})},
          ylabel={$\varepsilon$ (\si{V})}, xtick={0,0.01,...,0.081}, ytick={-4,-2,...,4},
          xticklabel style={/pgf/number format/.cd, fixed, fixed zerofill, precision=2}]
        \draw[hscink!70, line width=0.6pt] (axis cs:0,0) -- (axis cs:0.08,0);
        \addplot[hscseriesA, line width=1.3pt, mark=none, domain=0:0.08, samples=161]
          {4.0*sin(360*x/0.040)};
        \node[anchor=south, font=\footnotesize] at (axis cs:0.010,4.0) {$+4.0$};
        \node[anchor=north, font=\footnotesize] at (axis cs:0.030,-4.0) {$-4.0$};
      \end{axis}
    \end{tikzpicture}
    \caption*{Answer to \cref{pq:generator}(a): AC output for two turns.}
  \end{figure}

  \textbf{(b)} (i) With a split-ring commutator, the connections to the lamp swap every half
  turn, just when the emf is zero. So the output never reverses: every negative half becomes
  positive. It is DC that pulses between 0 and \qty{4.0}{V}, with a hump every
  \qty{0.020}{s}. (ii) The emf is induced in every turn, so it is proportional to the number
  of turns, $N$ ($\varepsilon = -N\dfrac{\Delta\Phi}{\Delta t}$). Twice as many turns double
  the peak emf to \qty{8.0}{V}. The coil turns at the same speed, so the period stays
  \qty{0.040}{s} (a frequency of \qty{25}{Hz}).

  \begin{modelanswer}
  \textbf{(c)} With the circuit open, an emf is induced but no current flows, so there is no magnetic
  force on the coil. With the lamp connected, the emf drives a current round the coil. The
  sides of the coil carry this current in the magnetic field, so they feel motor-effect
  forces ($F = lIB$). By Lenz's law, the induced current opposes the change that makes it,
  so these forces oppose the rotation. The person must do extra work against them.

  This extra work is where the lamp's energy comes from: the work done turning the handle
  becomes electrical energy, then light and heat in the lamp. Energy is conserved. If the
  forces helped the rotation instead, the generator would make energy from nothing.
  \end{modelanswer}
  \begin{markingcriteria}
    \criterion{(a) Sine curve that starts at zero, with a period of \qty{0.040}{s} (two full
      cycles in \qty{0.080}{s})}{1}
    \criterion{(a) Peak values of $\pm\qty{4.0}{V}$ shown}{1}
    \criterion{(b)(i) Split ring: the output is DC, with every half cycle on the same side
      (it never reverses)}{1}
    \criterion{(b)(ii) Twice the turns: peak \qty{8.0}{V}, period unchanged
      (\qty{0.040}{s})}{1}
    \criterion{(c) Lenz's law: the current in the coil gives forces that oppose the
      rotation (none when the circuit is open)}{1}
    \criterion{(c) Energy: the extra work done becomes the electrical energy in the lamp
      (energy is conserved)}{1}
  \end{markingcriteria}
\end{practiceans}

% ---- Q21 coil pulled through a field region
\begin{practiceans}
  \textbf{(a)}
  \given{$N = 25$; \ side $l = \qty{10}{cm} = \qty{0.10}{m}$; \ $v = \qty{2.0}{m.s^{-1}}$; \
    $B = \qty{0.40}{T}$; the plane of the coil is at \ang{90} to $B$, so $\theta = \ang{0}$}
  \find{the size of the emf while the coil enters the field}
  \formula{$\Phi = BA\cos\theta$ \ and \ $\varepsilon = -N\dfrac{\Delta\Phi}{\Delta t}$}
  \substitute{\eqlines{\Delta t &= \frac{\qty{0.10}{m}}{\qty{2.0}{m.s^{-1}}} = \qty{0.050}{s}
        \quad\text{(time to move fully into the field)}\\
      \Phi_\mathrm{final} &= \qty{0.40}{T} \times (\qty{0.10}{m})^2 \times \cos\ang{0}
        = \qty{4.0e-3}{Wb} \quad\text{(one turn)}\\
      \Delta\Phi &= \qty{4.0e-3}{Wb} - 0 = \qty{4.0e-3}{Wb}
        \quad\text{(the flux starts at zero)}\\
      \lvert\varepsilon\rvert &= 25 \times \frac{\qty{4.0e-3}{Wb}}{\qty{0.050}{s}}
        = \qty{2.0}{V}}}
  \answer{\final{$\varepsilon = \qty{2.0}{V}$}}
  \checkstep{motional emf (not on the formula sheet), one side of each turn cutting the
    field: $\varepsilon = NlvB = 25 \times 0.10 \times 2.0 \times 0.40 = \qty{2.0}{V}$
    \cmark.}

  \textbf{(b)} Follow the coil in four stages.
  \begin{itemize}
    \item \qtyrange{0}{0.050}{s}: the coil enters. The flux grows steadily, so the emf is
          $+\qty{2.0}{V}$.
    \item \qtyrange{0.050}{0.15}{s}: the whole coil is inside the uniform field. The flux
          stays at $BA$, so the emf is zero. (The right-hand side reaches the far edge of
          the field after moving \qty{30}{cm}:
          $t = \qty{0.30}{m} \div \qty{2.0}{m.s^{-1}} = \qty{0.15}{s}$.)
    \item \qtyrange{0.15}{0.20}{s}: the coil leaves. The flux falls steadily, so the emf is
          $-\qty{2.0}{V}$ (the opposite direction). The coil is fully out when it has moved
          \qty{40}{cm}: $t = \qty{0.20}{s}$.
    \item After \qty{0.20}{s}: no field, no flux, no emf.
  \end{itemize}
  \begin{figure}
    \centering
    \begin{tikzpicture}
      \begin{axis}[hsc, width=0.78\linewidth, height=0.30\linewidth,
          xmin=0, xmax=0.25, ymin=-3, ymax=3, xlabel={$t$ (\si{s})},
          ylabel={$\varepsilon$ (\si{V})}, xtick={0,0.05,...,0.251}, ytick={-3,-2,...,3},
          xticklabel style={/pgf/number format/.cd, fixed, fixed zerofill, precision=2}]
        \draw[hscink!70, line width=0.6pt] (axis cs:0,0) -- (axis cs:0.25,0);
        \addplot[hscseriesB, solid, line width=1.3pt, mark=none, const plot]
          coordinates {(0,2.0) (0.05,0) (0.15,-2.0) (0.20,0) (0.25,0)};
        \node[anchor=south, font=\footnotesize] at (axis cs:0.025,2.05) {$+2.0$};
        \node[anchor=north, font=\footnotesize] at (axis cs:0.175,-2.05) {$-2.0$};
      \end{axis}
    \end{tikzpicture}
    \caption*{Answer to \cref{pq:coilpull}(b): the emf in the coil.}
  \end{figure}

  \textbf{(c)}
  \given{$\varepsilon = \qty{2.0}{V}$; \ $R = \qty{0.40}{\ohm}$; \ $N = 25$; \
    $l = \qty{0.10}{m}$; \ $B = \qty{0.40}{T}$}
  \find{the force needed for a steady speed}
  \formula{$I = \dfrac{\varepsilon}{R}$ (from $V = IR$), \ then $F = lIB$ on the right-hand
    side of each of the $N$ turns}
  \substitute{\eqlines{I &= \frac{\qty{2.0}{V}}{\qty{0.40}{\ohm}} = \qty{5.0}{A}\\
      F &= N\,lIB = 25 \times \qty{0.10}{m} \times \qty{5.0}{A} \times \qty{0.40}{T}
        = \qty{5.0}{N}}}
  \answer{\final{$F = \qty{5.0}{N}$}, to the right (the same way as the motion).}
  \begin{modelanswer}
  As the coil enters, the flux into the page through it increases. By Lenz's law, the
  induced current opposes this increase. So it flows anticlockwise, making a field out of
  the page inside the coil. In the right-hand side this current flows up the page. The palm
  rule then gives a magnetic force to the left, against the motion. (The forces on the top
  and bottom sides are equal and opposite, so they cancel.) To keep a steady speed, the net
  force must be zero, so a pull of \qty{5.0}{N} to the right is needed.

  Energy: the pull does work at a rate $P = Fv = \qty{5.0}{N} \times \qty{2.0}{m.s^{-1}} =
  \qty{10}{W}$. This becomes electrical energy and then heat in the coil:
  $I^2R = (\qty{5.0}{A})^2 \times \qty{0.40}{\ohm} = \qty{10}{W}$. The two rates are equal,
  so energy is conserved. If the induced force helped the motion instead, the coil would
  speed up and make energy from nothing.
  \end{modelanswer}
  \begin{markingcriteria}
    \criterion{(a) Correct flux change (\qty{4.0e-3}{Wb}) and time (\qty{0.050}{s}), or uses
      $\varepsilon = NlvB$}{1}
    \criterion{(a) Correct emf, \qty{2.0}{V}}{1}
    \criterion{(b) Correct shape: a positive block, zero, a negative block, zero}{1}
    \criterion{(b) Correct values and times: $\pm\qty{2.0}{V}$; changes at 0.050, 0.15 and
      \qty{0.20}{s}}{1}
    \criterion{(c) Correct current, $I = \dfrac{\varepsilon}{R} = \qty{5.0}{A}$}{1}
    \criterion{(c) Correct force, $F = NlIB = \qty{5.0}{N}$ (all 25 turns counted)}{1}
    \criterion{(c) Explains: by Lenz's law the force on the induced current opposes the
      motion, AND the work done by the pull becomes electrical energy and heat}{1}
  \end{markingcriteria}
\end{practiceans}

% ---- Q22 parallel rods experiment (vary r)
\begin{practiceans}
  \textbf{(a)} $1/r$: change $r$ from millimetres to metres first (divide by 1000), then
  divide 1 by it. For example, $1 \div \qty{3.0e-3}{m} = \qty{333}{m^{-1}}$.
  $F$: change $\Delta m$ from grams to kilograms (divide by 1000), then multiply by $g$. For
  example, $F = \qty{0.256e-3}{kg} \times \qty{9.8}{m.s^{-2}} = \qty{2.51e-3}{N}$.
  \begin{center}
  \begin{tabular}{@{}lccccc@{}}
    \toprule
    $r$ (\si{mm}) & 3.0 & 4.0 & 5.0 & 6.0 & 8.0\\
    $1/r$ (\si{m^{-1}}) & 333 & 250 & 200 & 167 & 125\\
    $\Delta m$ (\si{g}) & 0.256 & 0.193 & 0.153 & 0.129 & 0.096\\
    $F$ ($\times 10^{-3}$ \si{N}) & 2.51 & 1.89 & 1.50 & 1.26 & 0.941\\
    \bottomrule
  \end{tabular}
  \end{center}

  \textbf{(b)} See the graph. The points lie close to a straight line through the origin.
  This shows that $F \propto 1/r$.
  \begin{figure}
    \centering
    \begin{tikzpicture}
      \begin{axis}[hsc, width=0.78\linewidth, height=0.34\linewidth,
          xmin=0, xmax=350, ymin=0, ymax=3.0,
          xlabel={$1/r$ (\si{m^{-1}})}, ylabel={$F$ ($\times 10^{-3}$ \si{N})},
          xtick={0,50,...,350}, ytick={0,0.5,...,3.0},
          yticklabel style={/pgf/number format/.cd, fixed, fixed zerofill, precision=1}]
        \addplot[hscseriesB, line width=1pt, domain=0:345, samples=2, mark=none]
          {0.0075377*x};
        \addplot[hscseriesA, only marks, mark=*, mark size=2pt]
          coordinates {(333.3,2.509) (250,1.891) (200,1.499) (166.7,1.264) (125,0.941)};
        \draw[hscmuted, densely dotted, line width=0.7pt] (axis cs:0,0) -- (axis cs:300,0)
          -- (axis cs:300,2.26);
        \node[anchor=west, font=\footnotesize, text=hscmuted] at (axis cs:302,1.15)
          {rise};
        \node[anchor=south, font=\footnotesize, text=hscmuted] at (axis cs:240,0.02)
          {run $= \qty{300}{m^{-1}}$};
        \node[anchor=south east, font=\footnotesize, text=hscseriesB] at (axis cs:150,1.4)
          {line of best fit};
      \end{axis}
    \end{tikzpicture}
    \caption*{Answer to \cref{pq:rods}(b): $F$ against $1/r$. The rise is
      \qty{2.26e-3}{N}.}
  \end{figure}

  \textbf{(c)}
  \given{two points on the line: $(0, 0)$ and $(\qty{300}{m^{-1}}, \qty{2.26e-3}{N})$; \
    $I = \qty{8.0}{A}$ in each rod; \ $l = \qty{0.60}{m}$}
  \find{the gradient, then $\mu_0$}
  \formula{$\dfrac{F}{l} = \dfrac{\mu_0}{2\pi}\dfrac{I_1 I_2}{r}$ with $I_1 = I_2 = I$ gives
    $F = \left(\dfrac{\mu_0 I^2 l}{2\pi}\right) \dfrac{1}{r}$. This is $y = kx$, so the
    gradient is $\dfrac{\mu_0 I^2 l}{2\pi}$ and
    $\mu_0 = \dfrac{2\pi \times \text{gradient}}{I^2 l}$.}
  \substitute{\eqlines{\text{gradient} &= \frac{\qty{2.26e-3}{N} - 0}{\qty{300}{m^{-1}} - 0}
        = \qty{7.533e-6}{N.m}\\
      \mu_0 &= \frac{2\pi \times \qty{7.533e-6}{N.m}}{(\qty{8.0}{A})^2 \times \qty{0.60}{m}}
        \quad\text{(substitute)}\\
      &= \frac{\qty{4.733e-5}{N.m}}{\qty{38.4}{A^2.m}} = \qty{1.233e-6}{N.A^{-2}}}}
  \answer{\final{$\mu_0 = \qty{1.23e-6}{N.A^{-2}}$} (values from about
    \qtyrange{1.21e-6}{1.25e-6}{N.A^{-2}}, from a good line, are fine)}
  \checkstep{the data-sheet value is $4\pi \times \num{e-7} = \qty{1.257e-6}{N.A^{-2}}$. The
    result is about 2\,\% lower: a good agreement \cmark. Read the power of ten from the axis
    label: 2.26 on the axis means \qty{2.26e-3}{N}. The gradient has the unit
    $\si{N} \div \si{m^{-1}} = \si{N.m}$.}

  \textbf{(d)} The formula is for infinitely long, straight wires. Near each end of a real
  rod, there is no more rod beyond the end. So the field of the other rod is weaker there,
  and the force is a little smaller than the formula predicts. When the rods are much longer
  than $r$ (here $l$ is 75 to 200 times $r$), the ends are only a tiny part of each rod.
  Then the formula is a good model. (This end effect makes the measured forces a little
  too small. It is one reason why the student's value of $\mu_0$ is a little lower than the
  data-sheet value.)
  \begin{markingcriteria}
    \criterion{(a) All values of $1/r$ and $F$ correct (mm changed to m, g changed to kg)}{1}
    \criterion{(b) Points plotted correctly}{1}
    \criterion{(b) Ruled straight line of best fit, through (or very close to) the
      origin}{1}
    \criterion{(c) Gradient from the line, about \qty{7.5e-6}{N.m}, with
      gradient $= \dfrac{\mu_0 I^2 l}{2\pi}$}{1}
    \criterion{(c) Correct $\mu_0$, about \qty{1.23e-6}{N.A^{-2}}}{1}
    \criterion{(d) The formula assumes very long (infinite) wires; the ends make the force
      smaller, so $l$ must be much bigger than $r$}{1}
  \end{markingcriteria}
\end{practiceans}

% ---- Q23 DC motor current graph
\begin{practiceans}
  \textbf{(a)}
  \given{$V = \qty{6.0}{V}$; at $t = 0$ the graph shows $I = \qty{4.0}{A}$. The coil is not
    yet turning, so there is no back emf.}
  \find{the resistance $R$ of the coil}
  \formula{$V = IR$, so $R = \dfrac{V}{I}$}
  \substitute{$R = \dfrac{\qty{6.0}{V}}{\qty{4.0}{A}} = \qty{1.5}{\ohm}$}
  \answer{\final{$R = \qty{1.5}{\ohm}$}}

  \textbf{(b)}
  \given{steady current with no load $I = \qty{0.40}{A}$ (from the graph);
    $R = \qty{1.5}{\ohm}$; $V = \qty{6.0}{V}$}
  \find{the back emf $\varepsilon_\mathrm{back}$}
  \formula{the back emf opposes the supply, so the net voltage drives the current:
    $V - \varepsilon_\mathrm{back} = IR$ (not on the formula sheet)}
  \substitute{$\varepsilon_\mathrm{back} = V - IR = \qty{6.0}{V} - \qty{0.40}{A} \times
    \qty{1.5}{\ohm} = \qty{6.0}{V} - \qty{0.60}{V} = \qty{5.4}{V}$}
  \answer{\final{$\varepsilon_\mathrm{back} = \qty{5.4}{V}$}}
  \checkstep{the back emf is less than the supply voltage, as it must be \cmark.}

  \begin{modelanswer}
  \textbf{(c)} At switch-on the coil is not moving, so there is no back emf. The current is
  limited only by the coil's resistance: $I = \dfrac{\qty{6.0}{V}}{\qty{1.5}{\ohm}} =
  \qty{4.0}{A}$. This large current gives a large torque ($\tau = nIAB$), so the motor speeds
  up.

  As the coil turns in the magnetic field, the flux through it changes, so an emf is induced
  in it (Faraday's law). By Lenz's law, this \textbf{back emf} opposes the supply voltage. The
  faster the coil turns, the bigger the back emf. So the net voltage
  ($V - \varepsilon_\mathrm{back}$) and the current fall as the motor speeds up. The graph
  falls steeply at first, then more slowly.

  The torque falls with the current. When the motor's torque just balances the friction
  torque, the speed stops rising. The back emf then stays constant (\qty{5.4}{V}), so the
  current stays at a small, steady \qty{0.40}{A}.

  \textbf{(d)} The load slows the motor. A slower coil has a smaller back emf (here
  $\qty{6.0}{V} - \qty{1.0}{A} \times \qty{1.5}{\ohm} = \qty{4.5}{V}$). So the net voltage
  is bigger and the current rises to \qty{1.0}{A}. The bigger current gives the bigger
  torque needed to turn the load. In energy terms, the motor now takes \qty{6.0}{W} from the
  supply instead of \qty{2.4}{W}, to do work on the load.
  \end{modelanswer}
  \begin{markingcriteria}
    \criterion{(a) $R = \qty{1.5}{\ohm}$, from the current at $t = 0$}{1}
    \criterion{(b) Back emf $= \qty{5.4}{V}$}{1}
    \criterion{(c) At switch-on there is no back emf, so the current is large (limited only by
      the resistance)}{1}
    \criterion{(c) As the speed rises, the back emf (induced, opposing the supply by Lenz's
      law) rises, so the current falls}{1}
    \criterion{(c) The current becomes steady when the motor torque balances the friction
      torque (constant speed, constant back emf)}{1}
    \criterion{(d) The load slows the motor, so the back emf falls and the current rises}{1}
  \end{markingcriteria}
\end{practiceans}

% ---- Q24 crossed fields
\begin{practiceans}
  \textbf{(a)}
  \given{$L = \qty{5.0}{cm} = \qty{0.050}{m}$; \ $d = \qty{2.0}{cm} = \qty{0.020}{m}$; \
    $V = \qty{120}{V}$; \ $v = \qty{3.0e7}{m.s^{-1}}$ (along the plates); \
    $q = \qty{1.602e-19}{C}$; \ $m_\mathrm{e} = \qty{9.109e-31}{kg}$}
  \find{the vertical deflection $s_y$ when the electron leaves the plates}
  \formula{$E = \dfrac{V}{d}$, \ $F = qE$, \ $a = \dfrac{F}{m}$; \ along the plates:
    $t = \dfrac{L}{v}$; \ across the plates: $s_y = u_y t + \frac{1}{2}at^2$ with $u_y = 0$}
  \substitute{\eqlines{E &= \frac{\qty{120}{V}}{\qty{0.020}{m}} = \qty{6000}{V.m^{-1}}\\
      a &= \frac{qE}{m_\mathrm{e}} = \frac{\qty{1.602e-19}{C} \times \qty{6000}{V.m^{-1}}}
        {\qty{9.109e-31}{kg}} = \qty{1.055e15}{m.s^{-2}}\\
      t &= \frac{\qty{0.050}{m}}{\qty{3.0e7}{m.s^{-1}}} = \qty{1.667e-9}{s}
        \quad\text{(no force along the plates)}\\
      s_y &= \tfrac{1}{2} \times \qty{1.055e15}{m.s^{-2}} \times (\qty{1.667e-9}{s})^2
        = \qty{1.466e-3}{m}}}
  \answer{\final{$s_y = \qty{1.5}{mm}$} upwards, towards the positive plate.}
  \checkstep{\qty{1.5}{mm} is less than the \qty{10}{mm} from the middle to a plate, so the
    electron leaves without hitting a plate \cmark. The electron is negative, so it is pulled
    towards the positive (top) plate \cmark. $a$ is about \num{e14} times $g$, so ignoring
    gravity is fine.}

  \textbf{(b)}
  \given{$E = \qty{6000}{V.m^{-1}}$; \ $v = \qty{3.0e7}{m.s^{-1}}$}
  \find{the size and direction of $B$}
  \formula{no deflection: the magnetic force balances the electric force, $qvB = qE$, so
    $B = \dfrac{E}{v}$ (not on the formula sheet)}
  \substitute{$B = \dfrac{\qty{6000}{V.m^{-1}}}{\qty{3.0e7}{m.s^{-1}}} = \qty{2.0e-4}{T}$}
  \answer{\final{$B = \qty{2.0e-4}{T}$}, \textbf{into the page}.}
  \checkstep{\emph{Direction.} The electric force on the electron is up, so the magnetic
    force must be down. Palm rule for a \emph{positive} charge moving right in a field into
    the page: fingers into the page, thumb to the right, palm faces up. The electron is
    negative, so its force is the opposite: down \cmark.}

  \textbf{(c)}
  \given{$m_\mathrm{e} = \qty{9.109e-31}{kg}$; \ $v = \qty{3.0e7}{m.s^{-1}}$; \
    $q = \qty{1.602e-19}{C}$; \ $B = \qty{2.0e-4}{T}$}
  \find{the radius $r$}
  \formula{$r = \dfrac{mv}{qB}$ (from $qvB = \dfrac{mv^2}{r}$; not on the formula sheet)}
  \substitute{$r = \dfrac{\qty{9.109e-31}{kg} \times \qty{3.0e7}{m.s^{-1}}}
    {\qty{1.602e-19}{C} \times \qty{2.0e-4}{T}}
    = \dfrac{\num{2.733e-23}}{\num{3.204e-23}} = \qty{0.853}{m}$}
  \answer{\final{$r = \qty{0.85}{m}$}}
  \begin{modelanswer}
  \textbf{Comparing the two motions.}
  \begin{itemize}
    \item \textbf{Force:} in the electric field the force ($qE$) is constant in size and
          direction (up). In the magnetic field the force ($qvB$) is constant in size, but
          it is always at \ang{90} to the velocity, so its direction keeps turning.
    \item \textbf{Path:} the electric field gives a parabola, like a projectile. The
          magnetic field gives an arc of a circle of radius \qty{0.85}{m} (uniform circular
          motion).
    \item \textbf{Direction:} the electric field bends the electron up, towards the positive
          plate. The magnetic field alone bends it down.
    \item \textbf{Speed:} the electric field does work on the electron, so its speed
          increases a little (to \qty{3.005e7}{m.s^{-1}}). The magnetic force is at \ang{90}
          to the motion, so it does no work: the speed stays \qty{3.0e7}{m.s^{-1}}.
  \end{itemize}
  (Over the \qty{5.0}{cm} plates, the arc bends the electron down by
  $r - \sqrt{r^2 - L^2} = \qty{1.5}{mm}$: the same size as in (a), but the other way. This
  fits with (b): the two forces are equal and opposite, so together they cancel.)
  \end{modelanswer}
  \begin{markingcriteria}
    \criterion{(a) Acceleration from $E = \dfrac{V}{d}$ and $F = qE$
      (\qty{1.06e15}{m.s^{-2}})}{1}
    \criterion{(a) Time between the plates, $t = \dfrac{L}{v} = \qty{1.67e-9}{s}$}{1}
    \criterion{(a) Deflection \qty{1.5}{mm}, towards the positive (top) plate}{1}
    \criterion{(b) $B = \dfrac{E}{v} = \qty{2.0e-4}{T}$}{1}
    \criterion{(b) Into the page, with a correct reason (the magnetic force must be down)}{1}
    \criterion{(c) Radius \qty{0.85}{m}}{1}
    \criterion{(c) Compares the paths: parabola (constant force) and circular arc (force at
      \ang{90} to the velocity)}{1}
    \criterion{(c) Compares speed or direction: the speed rises in $E$ but not in $B$, OR
      the electron bends the opposite way}{1}
  \end{markingcriteria}
\end{practiceans}

\end{practiceanswers}
